\section{O Algoritmo de Johnson e o Método dos Potenciais}

Na programação competitiva, o problema de Caminhos Mais Curtos entre Todos os Pares (APSP - \textit{All-Pairs Shortest Paths}) é frequentemente resolvido com o algoritmo de Floyd-Warshall em $\mathcal{O}(V^3)$. No entanto, para grafos esparsos, essa complexidade pode ser um gargalo. O Algoritmo de Johnson resolve o APSP em $\mathcal{O}(V E \log V)$, o que é muito mais eficiente em grafos onde $E \ll V^2$, mesmo na presença de arestas com pesos negativos.

\subsection{O Método dos Potenciais}

Se o grafo tiver apenas pesos não-negativos, poderíamos simplesmente executar o algoritmo de Dijkstra a partir de cada vértice, obtendo $\mathcal{O}(V E \log V)$. O problema surge com arestas negativas, que quebram a premissa gulosa do Dijkstra.

Para contornar isso, usamos o \textbf{Método dos Potenciais}. A ideia é atribuir um valor (potencial) $h(u)$ a cada vértice $u \in V$ e reponderar cada aresta $(u, v)$ de peso $w(u, v)$ com um novo peso $w'(u, v)$ definido como:
$$w'(u, v) = w(u, v) + h(u) - h(v)$$

Essa reponderação tem duas propriedades matemáticas fundamentais:
\begin{enumerate}
    \item \textbf{Preservação dos Caminhos Mais Curtos:} Para qualquer caminho $p = v_0 \rightarrow v_1 \rightarrow \dots \rightarrow v_k$, o novo peso total do caminho é:
    $$w'(p) = w(p) + h(v_0) - h(v_k)$$
    Como o termo $h(v_0) - h(v_k)$ depende apenas dos vértices de início e fim, todos os caminhos entre $v_0$ e $v_k$ sofrem exatamente a mesma alteração. Logo, o caminho mais curto no grafo original continua sendo o caminho mais curto no grafo reponderado.
    \item \textbf{Ciclos Negativos:} O peso de um ciclo não é alterado pela reponderação, pois o vértice de início e fim é o mesmo, o que faz com que $h(v_0) - h(v_k) = 0$.
\end{enumerate}

\subsection{O que torna um Potencial Válido?}

Nosso objetivo ao alterar os pesos é habilitar o uso do algoritmo de Dijkstra, que exige que todas as arestas possuam peso não-negativo para garantir a sua corretude. Assim, um potencial $h$ é considerado \textbf{válido} se, e somente se, para toda aresta $(u, v) \in E$, o novo peso não for negativo:
$$w'(u, v) \ge 0$$

Substituindo a definição de $w'$, isso significa que:
$$w(u, v) + h(u) - h(v) \ge 0$$
Isolando os potenciais, chegamos à restrição fundamental de um potencial válido:
$$h(v) - h(u) \le w(u, v)$$

Essa inequação tem exatamente o mesmo formato da desigualdade triangular que rege os caminhos mais curtos em grafos. 

Isso também nos prova matematicamente por que é impossível resolver o problema se houver um \textbf{ciclo de peso negativo}. Imagine um ciclo $C = v_1 \rightarrow v_2 \rightarrow \dots \rightarrow v_k \rightarrow v_1$. Se tentarmos encontrar um potencial válido e somarmos a inequação $h(v) - h(u) \le w(u, v)$ para todas as arestas desse ciclo, o lado esquerdo será exatamente $0$ (pois todo vértice aparece uma vez subtraindo e uma vez somando). O lado direito será a soma dos pesos do ciclo, $w(C)$. Teríamos então:
$$0 \le w(C)$$
Se $w(C)$ for um valor negativo, temos um absurdo matemático ($0 \le \text{negativo}$). Isso prova de forma elegante que, se o grafo possui um ciclo de peso negativo, é matematicamente impossível formular uma função potencial válida.

\subsection{O Algoritmo de Johnson}

O objetivo de Johnson é encontrar essa função potencial $h$ válida. Se garantirmos isso, podemos rodar o Dijkstra em segurança a partir de qualquer vértice. 

Os passos do algoritmo são:
\begin{enumerate}
    \item \textbf{Vértice Virtual:} Adicione um novo vértice $s$ ao grafo, com arestas de peso $0$ direcionadas para todos os outros vértices do grafo original.
    \item \textbf{Bellman-Ford:} Execute o algoritmo de Bellman-Ford a partir do vértice $s$. 
    \begin{itemize}
        \item Se o Bellman-Ford detectar um ciclo de peso negativo, o algoritmo para, pois o problema de caminho mais curto não tem solução.
        \item Caso contrário, defina o potencial de cada vértice $u$ como a distância mínima de $s$ até $u$, ou seja, $h(u) = \delta(s, u)$.
    \end{itemize}
    \item \textbf{Reponderação:} Pela desigualdade triangular do caminho mais curto, sabemos que $\delta(s, v) \le \delta(s, u) + w(u, v)$. Substituindo por $h$, temos $h(v) \le h(u) + w(u, v)$. Rearranjando a inequação, obtemos:
    $$w(u, v) + h(u) - h(v) \ge 0$$
    Logo, definimos $w'(u, v) = w(u, v) + h(u) - h(v)$ e agora garantimos que todas as arestas são não-negativas.
    \item \textbf{Dijkstra Múltiplo:} Remova o vértice virtual $s$. Para cada vértice $u \in V$, execute o algoritmo de Dijkstra no grafo reponderado para encontrar as menores distâncias $d'(u, v)$.
    \item \textbf{Reajuste das Distâncias:} A distância real no grafo original é recuperada revertendo a fórmula do potencial:
    $$d(u, v) = d'(u, v) - h(u) + h(v)$$
\end{enumerate}

Na prática das maratonas de programação, o Algoritmo de Johnson é a ferramenta ideal quando você se depara com um grafo grande e esparso, contendo arestas de peso negativo, e que exigiria múltiplas consultas, situação em que o algoritmo de Floyd-Warshall resultaria inevitavelmente em \textit{Time Limit Exceeded} (TLE).
